% Appendix A: TIDAL pipeline architecture (~1800 w + diagram).
% Per genre_conventions_software.md §(h). The audience is a reader who
% wants to understand the overall system before examining any subsystem.
% Model: HiGGS appendix (2206.00658) — primary-data style.

\section{TIDAL pipeline architecture}%
\label{Architecture}

\paragraph*{Overview}
% TODO: 3–4 sentences précis of the whole pipeline — inputs, outputs,
% what is automated. Open the appendix; do not repeat the §3 summary.

% Figure A.1 — TIDAL architecture diagram.
\begin{figure}[t]
  \centering
  % TODO: \documentclass[border=10pt]{standalone}
\usepackage{tidal-tikz-styles}

\begin{document}
\begin{tikzpicture}[
    node distance=0.4cm
]

% ===================================================================
% ROW 1: Main pipeline stages (left to right)
% ===================================================================

% --- Input ---
\node[data] (toml) {\texttt{theory.toml}};

% --- Derive ---
\node[stage=Pink, right=0.6cm of toml] (derive)
    {Symbolic derivation\\[2pt]\scriptsize Wolfram/xAct};

% --- JSON (narrow waist) ---
\node[data, right=0.6cm of derive, minimum width=1.8cm,
      draw=BlueLine, fill=BlueFill, fill opacity=0.35,
      font=\footnotesize\bfseries] (json)
    {\texttt{spec.json}};

% --- Simulate ---
\node[stage=Purple, right=0.6cm of json] (sim)
    {Numerical integration};

% --- Snapshots ---
\node[data, right=0.6cm of sim] (snaps)
    {Field\\[-1pt]\scriptsize snapshots};

% --- Measure ---
\node[stage=Green, right=0.6cm of snaps] (meas)
    {Post-processing};

% --- Results ---
\node[data, right=0.6cm of meas] (results)
    {Measurement\\[-1pt]\scriptsize results};

% ===================================================================
% ARROWS (main flow)
% ===================================================================
\draw[arr] (toml) -- (derive);
\draw[arr] (derive) -- (json);
\draw[arr] (json) -- (sim);
\draw[arr] (sim) -- (snaps);
\draw[arr] (snaps) -- (meas);
\draw[arr] (meas) -- (results);


% ===================================================================
% BRACE: Symbolic vs Numerical separation at JSON
% ===================================================================
\draw[decorate, decoration={brace, amplitude=6pt, raise=2pt},
      PinkLine, thick]
    (toml.west |- derive.north) -- (derive.north east)
    node[midway, above=10pt, font=\scriptsize\bfseries, text=PinkLine]
    {Symbolic (CAS)};

\draw[decorate, decoration={brace, amplitude=6pt, raise=2pt},
      PurpleLine, thick]
    (sim.north west) -- (meas.north east)
    node[midway, above=10pt, font=\scriptsize\bfseries, text=PurpleLine]
    {Numerical (PDE)};

% JSON label
\node[above=0.25cm of json, font=\scriptsize\bfseries, text=BlueLine]
    {portable contract};

\end{tikzpicture}
\end{document}
 (TikZ)
  % TODO caption: "TIDAL pipeline. The symbolic stage (\xAct/\Mathematica)
  % ingests a TOML Lagrangian, derives the Euler–Lagrange equations,
  % decomposes by component, and exports a JSON specification. The
  % numerical stage (\Python) dispatches to one of five solver backends
  % based on equation properties, producing simulation data and
  % measurements."
  \caption{TODO: pipeline diagram (symbolic --> JSON --> numerical).}
  \label{ArchitectureDiagram}
\end{figure}

\paragraph*{Symbolic stage}
% TODO: 4–6 sentences. \xAct/\xPert, Euler–Lagrange derivation, component
% decomposition. User-perspective only — no internal routine descriptions.
% This paragraph is a high-level summary; the detailed treatment lives in
% \cref{SymbolicComputation}.

\paragraph*{JSON handoff}
% TODO: 4–6 sentences. JSON schema fields present (equations, fields,
% canonical block, hamiltonian terms, etc.); reference json_schema.tex
% if appropriate.

\paragraph*{Solver dispatch}
% TODO: 4–6 sentences. Auto-selection logic — modal first if conditions
% are met, otherwise IDA/CVODE/leapfrog/scipy fallback. The full backend
% description lives in \cref{NumericalEvolution}.

\paragraph*{Output}
% TODO: 4–6 sentences. Simulation data files, measurement types
% (energy, conversion, mixing, spectrum, etc.).

\paragraph*{Dependencies}
% TODO: explicit citations to each package: \xAct, \xPert, xTras,
% SUNDIALS/scikit-sundae, \NumPy, \SciPy, \Mathematica. Cite via
% \cite{} for each — bib audit must verify these keys exist.
